Avec l’évolution continue des technologies numériques, les plateformes en ligne occupent aujourd’hui une place centrale dans la mise en relation des utilisateurs et la circulation de l’information. Elles facilitent les échanges, automatisent certains processus et offrent un accès rapide à divers services, quel que soit le domaine d’application.

Dans le secteur des animaux domestiques, ces solutions jouent un rôle particulièrement important. Elles permettent de répondre à des besoins variés tels que l’adoption, l’accouplement, la recherche d’animaux perdus ou encore l’accès à des services spécialisés comme les soins, la garde ou la vente d’accessoires. Ces fonctionnalités contribuent à améliorer la gestion et le bien-être des animaux, tout en simplifiant les interactions entre les utilisateurs.

Cependant, malgré la diversité des outils existants, ces solutions restent souvent fragmentées et peu centralisées. Les utilisateurs doivent alors naviguer entre plusieurs plateformes, ce qui complique l’expérience, réduit l’efficacité et limite la visibilité des informations disponibles. Cette dispersion rend donc nécessaire la conception d’une solution unifiée et mieux structurée.

Dans ce contexte, notre projet de fin d’études vise à concevoir et développer une plateforme dédiée aux animaux domestiques, regroupant plusieurs modules complémentaires. Cette plateforme permet de centraliser la gestion des annonces, de faciliter la mise en relation entre utilisateurs et de proposer divers services adaptés aux besoins du domaine animalier, le tout dans un environnement intuitif et performant.

Le présent rapport décrit les différentes étapes de réalisation de ce projet, depuis l’analyse des besoins jusqu’à la mise en œuvre des fonctionnalités. Il est structuré en quatre chapitres comme suit :

\begin{itemize}

\item \textbf{Chapitre 1 : Contexte général du projet} \\
Dans ce chapitre, nous allons présenter le cadre général du projet, la problématique ainsi que les objectifs visés.

\item \textbf{Chapitre 2 : Étude des besoins et planification du projet} \\
Dans ce chapitre, nous allons analyser les besoins fonctionnels et non fonctionnels, identifier les acteurs du système et présenter la planification adoptée.

\item \textbf{Chapitre 3 : AniHub Partie 1} \\
Dans ce chapitre, nous allons décrire la première release du projet AniHub, qui inclut le module d’authentification ainsi que les fonctionnalités d’administration, notamment AniMatch et AniFound.

\item \textbf{Chapitre 4 : AniHub Partie 2} \\
Dans ce chapitre, nous allons présenter la deuxième release du projet AniHub, comprenant le module AniMarket, le module AniCare, ainsi que les aspects de communication et de déploiement du système.

\end{itemize}